\chapter*{Prologue}
\addcontentsline{toc}{chapter}{Prologue}

The prologue serves as a thematic or narrative introduction to your dissertation. It can establish the emotional, intellectual, or conceptual landscape that frames your work.

\section*{Purpose and Function}

A prologue might:
\begin{itemize}
\item Present a compelling case study or example that illustrates your research problem
\item Tell a story or anecdote that motivates your inquiry
\item Establish key tensions, paradoxes, or questions that drive your research
\item Create a narrative arc that connects your introduction to your main argument
\item Engage readers on a more human or intuitive level before diving into scholarly analysis
\end{itemize}

\section*{Distinguishing Features}

While the introduction is structured and analytical, the prologue can employ more narrative or rhetorical strategies. It sets the stage and draws readers in, preparing them for the intellectual journey ahead.

\section*{Integration with Main Text}

Your prologue should connect meaningfully to Chapter 1 (Introduction), preparing readers for the formal presentation of your research questions, objectives, and methodology.

\newpage